% Methods section for "Graphical Processing Units in Science".
% Every quantity below is read from the pipeline's own output tables
% (registry v0.3.0; extraction 2026-07/08). Regenerate with the notebooks in
% notebooks/ and notebooks/arxiv/ if the pipeline is re-run.
%
% KNOWN INCOMPLETENESS (see \S3.5): S2ORC passage shards 6, 27 and 28 --
% 75,000 of 877,101 candidate passages (8.6\%) -- have no mentions table, and
% the missing passages are skewed toward recent publication years. All S2ORC
% numerators below and in Figs. 3, 5, 7, 9, 11, 12, 14, 15 are therefore
% computed on 91.4\% of candidate passages. Running
% `python -m accelscan.infer --shards 6,27,28` and rebuilding stages 3.5--4
% closes the gap and will change every S2ORC count.

\subsection{Corpora}

We measure reported accelerator use in two full-text corpora, processed by
identical code and reported separately.

\textbf{S2ORC.} We use the Semantic Scholar Open Research Corpus (S2ORC v2)
full-text release, a 214-shard gzipped JSONL snapshot retrieved 2026-03-13.
Each record carries a machine-extracted body together with character-offset
annotations for paragraphs and section headers; the reference list is a
separate top-level object, so bibliographies are excluded by construction. Paper
metadata (publication year, fields of study, citation counts) comes from the
matching Semantic Scholar metadata snapshot (2026-02-10). The corpus contains
$12{,}408{,}584$ distinct papers with full text, of which $11{,}790{,}600$
fall in the $2005$--$2025$ analysis window.

\textbf{arXiv.} We additionally process the arXiv \emph{author-submitted
\LaTeX{} source} bulk release (\texttt{s3://arxiv/src}), $12{,}830$
requester-pays tar archives covering 1991--2026 and containing $3{,}091{,}651$
submissions. Source is streamed rather than mirrored: archives are read
forward-only in-region on EC2 (\texttt{us-east-1}), where S3-to-instance
transfer is free, and only derived parquet is written back to institutional
storage. Of all submissions, $249{,}638$ ($8.1\%$) are PDF-only and carry no
source, $21{,}512$ contain no \TeX{} member, $4{,}358$ are PostScript,
$2{,}131$ exceed the size ceiling, and $40$ are corrupt, non-gzip, or exceeded
the per-paper time budget; submission shapes are classified by magic bytes
rather than file extension. $2{,}813{,}972$ submissions ($91.0\%$) yield at
least one usable paragraph. The resulting population is $2{,}830{,}540$ papers,
$2{,}541{,}291$ of them in the $2005$--$2026$ window. arXiv adds text fidelity
(author \LaTeX{} rather than PDF extraction), coverage of the current month,
and arXiv's own category taxonomy; it is a self-selected preprint corpus skewed
toward physics, mathematics and computer science, so we compare
\emph{within-corpus trends} only, never cross-corpus levels.

Identity is unified by a corpus-namespaced key \texttt{paper\_id} --- the
Semantic Scholar \texttt{corpusid} for S2ORC, \texttt{arxiv:2301.01234} for
arXiv --- and a passage key \texttt{paper\_id:para\_idx}. arXiv
publication year is recovered from the identifier itself (every arXiv id
encodes \texttt{YYMM}, with $\mathrm{YY}\ge 91$ denoting the 1900s), which
agrees with the v1 submission date recorded in arXiv's own metadata snapshot
for $99.71\%$ of its $3{,}113{,}330$ records.

\subsection{Paragraph segmentation}

Both corpora are reduced to the same unit --- an ordered list of paragraphs, each
with the nearest preceding section heading --- so that all matching, gating and
passage assembly downstream is performed by one shared implementation.
Paragraphs shorter than 30 characters are dropped, and paragraph indices are
assigned before that filter so that \texttt{passage\_id} means the same thing in
both corpora.

For S2ORC we use the supplied paragraph and section-header offsets, falling back
to blank-line splitting for the minority of records with null annotations.

For arXiv we search the \LaTeX{} source \emph{as written}, without de-\TeX{}ing.
Hardware names are ASCII literals that appear verbatim in source, and two
parsing implementations (a hand-rolled converter and \texttt{pylatexenc}) were
evaluated and rejected: their measured benefit over raw search was confined to
\emph{exclusions}, while the risk --- non-termination on the routinely malformed
markup of a 2.6M-paper corpus, and era-correlated conversion failure that would
manufacture the very time trend we measure --- was not acceptable. All text
members of a submission are concatenated (longest first) rather than resolving
\verb|\input|, which is both simpler and higher-recall, since arXiv submissions
are typically multi-file and the methods section is often its own file. The
preamble is removed per file at \verb|\begin{document}|
($99.1\%$ of papers), paragraphs are split on blank lines (\TeX{}'s own rule)
with over-long paragraphs re-split at sentence boundaries to fit the passage
budget, section headings are captured from
\verb|\chapter|/\verb|\section|/\verb|\subsection|/\verb|\paragraph|, and
purely numeric inline math is unwrapped so that ``\verb|$8$ V100s|'' retains its
device count. Decoding falls back utf-8 $\rightarrow$ cp1252 $\rightarrow$
latin-1 ($93.9\%$, $5.8\%$, $0.3\%$ of papers respectively). Papers have a
median of 77 paragraphs.

Bibliographies are excluded in three layers: \texttt{.bbl} members are never
read; the document is truncated at the first
\verb|\begin{thebibliography}|, \verb|\bibliography{|, or
\verb|\printbibliography| occurring past the halfway point of the document (so
that a preamble \verb|\bibliographystyle| cannot decapitate a paper); and a
final filter discards reference-shaped paragraphs (numbered or \verb|\bibitem|
openers, or short paragraphs carrying two or more bibliographic signals). This
matters quantitatively: reference titles routinely contain ``GPU-accelerated
\ldots'', and $91.9\%$ of arXiv papers reach the truncation rule, with
$5{,}186{,}759$ reference-shaped paragraphs additionally filtered.

The accepted cost of raw-source search is that figure captions, spec tables,
code listings and commented-out text (\verb|% we used to use a K80|) are
searched. The bias is upward, known in sign, and concentrated in the
any-mention series; the usage-context stage (\S\,``Mention versus reported use'') is the primary defence,
since a caption or a related-work sentence is labelled as not
\emph{used-in-this-work}. This makes arXiv any-mention prevalence biased upward
relative to S2ORC, which obtains these exclusions free from PDF extraction.

\subsection{Hardware registry}

Matching is driven by a versioned registry (\texttt{registry/}, semantically
versioned; every output table is namespaced by the registry version that
produced it). Model entries are \emph{generated}, not authored: a scraper
(\texttt{scripts/build\_registry.py}) parses the specification tables of the
Wikipedia lists of NVIDIA and AMD graphics processing units, the AMD Instinct
page, the Tensor Processing Unit page, the ``Comparison of M-series processors''
table of the Apple silicon page, the Xeon Phi coprocessor and many-core model
tables, and the desktop and workstation tables of the Intel Arc page, collapsing
per-variant rows (memory size, form factor, GPU-core bin, edition) into canonical
models and merging duplicate entries
across pages by taking the earliest launch date, the union of aliases, and the
element-wise maximum of specifications. A second generator
(\texttt{scripts/build\_epoch\_registry.py}) adds the accelerators those pages
omit, from Epoch AI's machine-learning-hardware table
(\url{https://epoch.ai/data/machine-learning-hardware}, 176 devices, 2008--2026):
Chinese domestic silicon (Huawei Ascend, Cambricon MLU, Baidu Kunlun, Biren,
MetaX, Iluvatar), first-party cloud parts (Meta MTIA, Microsoft Maia, AWS
Trainium and Inferentia generations), export-control SKUs (A800, HGX H20,
MI308X), and datacenter GPUs missing from the Wikipedia lists, of which the
consequential one is the T4. It emits only devices the registry does not already
resolve to a model entry, and Epoch's FLOPS convention was verified to match ours
before use --- dense and non-sparse (H100 SXM tensor FP16 $=989.4$, not $1979$,
TFLOPS), reproducing the scraped Wikipedia figures for shared parts. Only rows with a parseable launch date
$\ge 2000$ are kept; mobile, integrated, console and chipset sections are
skipped. No specification number is ever taken from model memory: each entry
records a \texttt{spec\_source} naming the source page and table heading.
Hand curation (\texttt{registry/hardware.yaml}) supplies only what cannot be
scraped --- context-gate vocabularies, generic and architecture entries, and
overrides, which replace generated entries by id.

Registry v0.3.0 contains $1{,}236$ entries: $1{,}217$ specific models, 10
NVIDIA architectures (Fermi, Kepler, \ldots), and 9 generic terms (\texttt{GPU},
\texttt{CUDA}, \texttt{TPU}, \ldots, first-class entries so that generic-only
reporting is measurable), spanning 22 manufacturers. The largest are NVIDIA 568,
AMD 546, Intel 43, Apple 18 and Google 11, then a long tail of eight or fewer
entries each (Huawei, Cambricon, Baidu, Amazon, Meta, Microsoft, MetaX, Biren,
Iluvatar, \ldots). By market segment: consumer 929, workstation 136,
datacenter 130, cloud 22. Every specific model is generated; the hand-maintained
file holds only the gate vocabularies, the generic terms, two vendor brand
triggers and the architecture words, none of which any hardware table supplies.
Two classes are deliberately descoped and carry no model entry: embedded
inference modules (the NVIDIA Jetson line) and wafer-scale/IPU parts (Cerebras,
Graphcore), for which no source we parse publishes per-device specifications.
Their mentions are bucketed as out-of-scope in the unresolved audit rather than
counted as registry gaps, and they never enter a model-level series. Architecture
entries exist for NVIDIA only, whose architecture names are routinely used as
standalone hardware references (``trained on Hopper''); AMD's GCN/RDNA/CDNA are
carried as a metadata column on the model entries instead, since an
architecture-only AMD mention occurs in 2 papers per corpus.
Apple silicon is carried per chip variant (M1, M1 Pro, M1 Max, M1 Ultra, \ldots,
M5 Max) rather than per generation, because within a generation the variants
differ by up to $8\times$ in GPU throughput and $16\times$ in addressable memory.
Release dates span 2000--2026 and are populated for all $1{,}217$ models, since
the frontier-lag and half-life estimands depend on them. Per-device
specification coverage across models is: memory $98.0\%$, FP32 $64.3\%$, FP64
$45.9\%$, dense FP16 tensor throughput $7.8\%$ (defined only for post-Volta
parts), TDP $70.9\%$. Where a family is named without a model number --- ``an
Intel Xeon Phi'', ``AWS Trainium'' --- the registry carries a family entry whose
specifications are null, since an unnumbered mention does not say which
generation ran. Missing specifications remain null and are excluded
per-axis rather than zero-filled.

Following the FLOPS convention required for cross-generation comparability, we
record \emph{dense, non-sparse} peak throughput at boost clock, per device:
sparsity-doubled marketing figures are excluded, TF32 tensor throughput is never
recorded in the FP32 column, and only a matrix-engine column may feed the tensor
axis (Intel Arc's plain half-precision column is vector FP16 and is not read as
tensor throughput). Xeon Phi parts report only the peak FP64 figure Intel
published, so their FP32 and tensor axes are null rather than back-computed.

The two corpora were scanned under adjacent registry versions --- S2ORC under
v0.1.0 ($1{,}127$ entries, $1{,}838$ aliases; sources retrieved 2026-07-13) and
arXiv under v0.2.0 ($1{,}139$ entries, $1{,}862$ aliases; retrieved
2026-07-28). The bump changed no existing alias; it added per-device
specifications to existing entries and 12 new models, all of them the GeForce
RTX 40 series, which a malformed table cell had caused the scraper to drop. We
bound the resulting S2ORC recall loss by measurement on arXiv, which was scanned
with those aliases present: $14{,}728$ arXiv papers match an RTX 40 alias, but
$98.1\%$ of them are also triggered by another registry term in the same
passage, so only $282$ papers ($0.09\%$ of candidate papers) and $1{,}257$
passages ($0.15\%$) would fail to be generated without them. Because the LLM
reads whichever passages are generated and canonicalization runs against the
current registry, RTX 40 attribution itself is unaffected. Both capacity tables use the
current registry's specifications (v0.3.0, which additionally splits Apple
silicon into per-variant entries); specification additions do not change matching
and so require re-running only stages 3.5 and later, never the scan or the LLM.

\subsection{Candidate passage generation}

Stage 1 is a single CPU pass over each corpus emitting two products: an
\emph{inventory} row per paper (the population), and a \emph{candidate} row per
matched paragraph.

Matching is three-layered. (i) An Aho--Corasick automaton over the lowercase
literal cores of every alias in the extraction-time registry ($1{,}862$ for
arXiv, $1{,}838$ for S2ORC) prefilters paragraphs; a hit must fall on
word boundaries in the haystack. The overwhelming majority of paragraphs stop
here, which is what makes a full-corpus sweep affordable. (ii) For each
surviving core, a confirming per-alias regular expression is applied, wrapped in
hardware-safe boundaries \verb|(?<!\w)(?:...)s?(?!\w)| so that \texttt{V100}
does not match \texttt{CV100x} or \texttt{V1000}; literal spaces compile to
\verb|[\s-]*| so \texttt{RTX 3090}, \texttt{RTX-3090} and \texttt{RTX3090} are
one alias. Case sensitivity is derived from the pattern: all-uppercase aliases
(\texttt{GPU}, \texttt{V100}, \texttt{K80}) match case-sensitively,
mixed-case ones (\texttt{Tesla V100}) case-insensitively. Overlapping matches
are resolved by longest span, so \texttt{RTX 3090 Ti} beats \texttt{RTX 3090}
without hand-authored lookaheads. (iii) Ambiguous aliases --- short codes whose
strings collide with non-hardware usage (\texttt{K80} the antibody, \texttt{M2}
the money supply, architecture names that are ordinary words) --- are marked
tier~B ($236$ of the arXiv scan's $1{,}862$ aliases) and count only if a context term from the
relevant gate vocabulary (\texttt{gpu}: 21 terms; \texttt{tpu}: 5;
\texttt{apple}: 8) occurs within $\pm 250$ characters of the match, taken over
the paragraph and its immediate neighbours. A paper-level rescue then un-gates
tier~B matches in papers carrying at least one ungated model hit anywhere.
Gating errs permissive by design: it exists to bound candidate volume, not to be
precise, because precision is the LLM's job.

Each matched paragraph yields one candidate passage: the paragraph plus its
immediate neighbours, capped at $2{,}500$ characters ($\approx 600$ tokens). The
window is centred on the matched span rather than taken from the head of the
paragraph, so that the trigger is always inside the passage it triggered; every
recorded match offset is verified to index the stored passage text. Passages per
paper are capped at 20 for papers with at least one model-specific hit and 5 for
generic-only papers (\texttt{GPU}/\texttt{CUDA} with no model named), with a
\texttt{passages\_truncated} flag preserving auditability; generic-only papers
are retained rather than discarded because device counts and usage context are
measurable without a model name. Auxiliary device-count and memory expressions
adjacent to a hit are recorded as \texttt{matched\_surfaces} and passed to the
LLM as hints, but never trigger a passage on their own.

Stage 1 yields $877{,}101$ candidate passages from $399{,}857$ S2ORC papers and
$863{,}815$ from $312{,}480$ arXiv papers. These are repacked into fixed-size
shards of $25{,}000$ passages (36 and 35 shards respectively) so that GPU tasks
have uniform runtimes.

\subsection{LLM extraction and validation}

Stage 2 submits every candidate passage to a locally hosted open-weight model,
which decides whether each registry hit is genuinely an accelerator mention and,
if so, labels it. We use \textbf{Qwen3-14B} in bfloat16 served by vLLM's offline
batch interface, with a $4{,}096$-token context, at most $1{,}024$ output
tokens, greedy decoding ($\text{temperature}=0$), reasoning/``thinking'' mode
disabled, and prefix caching enabled (the system prompt is shared across all
passages in a batch).

Output is constrained by \textbf{guided JSON decoding} against a schema compiled
from a Pydantic model, so the model cannot emit unparseable output or
out-of-vocabulary labels. Each mention object carries: \texttt{model\_raw}
(verbatim surface form), \texttt{model\_normalized} (canonical marketing name,
or null when only a generic term appears), \texttt{manufacturer} (11-way enum),
\texttt{accelerator\_subtype} (11-way enum spanning datacenter/consumer/
workstation/mobile GPU, TPU, FPGA, NPU-ASIC, manycore, the two generic classes,
and \texttt{not-an-accelerator}), \texttt{device\_count} with
\texttt{device\_count\_basis} $\in$ \{explicit, inferred\}, \texttt{memory\_gb}
(per-device, when stated), \texttt{usage\_context} (below), and
\texttt{evidence\_quote}, a short verbatim span supporting the labels that
permits post-hoc hallucination auditing. The passage-level object is a list of
mentions, so an empty list is the model's explicit way of rejecting a registry
false positive; \texttt{not-an-accelerator} is the second off-ramp.

The prompt is a fixed system message (task definition, per-field rules,
definitions of the four usage contexts, and four worked examples --- a
multi-GPU methods sentence, a passage mixing related-work TPUs with a
used-in-this-work consumer GPU, an architecture studied as an object with a card
used for validation, and the ``K80 cells were incubated with the antibody''
negative) followed by a per-passage user message containing the section header,
the registry hint surfaces, and the passage text. Prompt text is versioned
(\texttt{PROMPT\_VERSION}) and stamped, with the registry version, model name
and code commit, into every output row; output tables are namespaced by
(registry version, prompt version, model), so variants can never be silently
mixed.

Exactly one row is written per passage even when nothing is extracted, with
\texttt{status} $\in$ \{\texttt{ok}, \texttt{no\_mention}, \texttt{truncated},
\texttt{json\_error}\}, and the raw generation is retained so that reprocessing
never requires re-inference. Structured decoding is highly reliable in practice:
of $1{,}010{,}417$ S2ORC output rows, $955{,}182$ ($94.5\%$) are \texttt{ok},
$55{,}141$ ($5.5\%$) report no accelerator mention (i.e., an explicitly rejected
registry candidate), 92 hit the output-token limit and 2 failed to parse; the
arXiv figures are $972{,}174$ ($90.8\%$), $97{,}608$ ($9.1\%$), 495 and 143 of
$1{,}070{,}420$ rows. The higher arXiv rejection rate is the expected
consequence of searching raw source.

Throughput is roughly one GPU-hour per $25{,}000$-passage shard
($\approx 680$ output tokens/s), i.e.\ $\approx 35$ GPU-hours per corpus. Jobs
run as SLURM arrays on a preemptible partition; the output shard is the unit of
work and a zero-byte \texttt{.done} marker the unit of completion, so preemption
costs at most one shard and reruns skip finished work.
% NOTE: 33 of 36 S2ORC shards completed; see the header comment.

\subsection{Normalization}

Stage 3 resolves each distinct \texttt{model\_normalized} string to a canonical
registry id by running the registry matcher over the string itself and keeping
the longest match, breaking ties toward model over architecture over generic.
Context gating is deliberately retained at this step, so that a bare
``Fermi'' in a hardware context becomes the architecture while ``Fermi
Gamma-ray Space Telescope'' does not become a GPU. Because canonicalization is
a table join over distinct strings, adding registry entries requires re-running
only this stage, never the LLM. \texttt{manufacturer} is taken directly from the
model's enum output; the registry is used only for model identity.

Strings that fail to canonicalize are classified rather than discarded, so that
the extracted-but-unregistered share is explicit. Among used-in-this-work
mentions naming a model, $94.9\%$ (S2ORC) and $97.0\%$ (arXiv) resolve to a
registry entry. The largest unresolved buckets are deliberately out-of-scope
accelerator families --- Xilinx FPGAs ($2.7\%$ / $1.3\%$ of named mentions),
Intel FPGAs, and CPUs --- which are excluded because the study is GPU- and
TPU-centric and a nameplate FLOPS figure for an FPGA is bitstream-dependent and
therefore meaningless. Genuinely unknown strings are $1.4\%$ (S2ORC) and
$1.0\%$ (arXiv).

\subsection{Mention versus reported use}

A textual mention does not establish that a paper used the device. Every mention
is therefore assigned exactly one usage context: \emph{used-in-this-work} (the
authors ran computations for this paper on it), \emph{comparison-or-related-work}
(prior work, another paper's baseline, or a specification comparison the authors
did not run), \emph{object-of-study} (the hardware is what the paper
benchmarks, simulates, or designs against), or \emph{speculative-future}
(announced or hypothetical hardware). Paper-level context is the
highest-priority context across the paper's mentions.

The distinction is quantitatively large and not a formality. Of $955{,}182$
confirmed S2ORC mentions, $501{,}451$ ($52.5\%$) are used-in-this-work,
$341{,}555$ ($35.8\%$) comparison-or-related-work, $93{,}878$ ($9.8\%$)
object-of-study, and $18{,}298$ ($1.9\%$) speculative; the arXiv shares are
$54.1\%$, $35.8\%$, $8.7\%$ and $1.5\%$ of $972{,}174$ mentions. Treating any
mention as use would inflate the mention count by $1.9\times$ and the count of
accelerator papers by $1.3\times$ ($339{,}601$ versus $258{,}312$ in S2ORC), and
the two series diverge over time. All primary analyses condition on
used-in-this-work; mention-imputed variants are reported as robustness. The
usage-context stage also absorbs much of the raw-source bias described above: it is what allows a corpus in which captions and comments are searched
to yield a defensible \emph{use} series.

We further distinguish \emph{model-specific} from \emph{generic-only} reporting.
A paper reporting use but naming no model carries no vendor, vintage, or
capacity information, and the generic-only share falls sharply over the study
period; it is therefore plotted before any composition or share figure, since a
falling unknown share can mechanically inflate apparent vendor growth.

\subsection{Reported compute capacity}

Stage 3.5 joins per-device registry specifications to canonicalized
used-in-this-work mentions. For paper $i$ and used model $g$ with device count
$\hat c_{ig}$,
\begin{align}
F_i^{(p)} &= \sum_{g \in \mathrm{used}(i)} \hat c_{ig}\,\mathrm{flops}_p(g),
  \qquad p \in \{\text{FP32}, \text{FP64}, \text{FP16 tensor}\},\\
M_i &= \sum_{g} \hat c_{ig}\,\mathrm{vram}(g), \qquad
W_i = \sum_{g} \hat c_{ig}\,\mathrm{tdp}(g),
\end{align}
with $C_i = \max_g \hat c_{ig}$ the scale-out measure. Mentions are first
deduplicated to one row per (paper, model) taking the maximum reported device
count, so a model repeated across passages is not double-counted. A per-device
memory figure stated in text resolves configuration variants (e.g.\ A100 40 vs.\
80\,GB) in preference to the registry default; stated memory is available for
$11.2\%$ (S2ORC) and $12.5\%$ (arXiv) of confirmed mentions, which also supplies
a validation subset for the specification-imputed series.

Two spec conventions are worth stating because they are visible in the results.
Integrated Apple GPUs have no FP64 figure (Metal exposes no double precision) and
no dense FP16 tensor figure (the Neural Engine's integer TOPS is a different
engine and a different unit), so they are null on those axes and excluded from
them rather than imputed; and their memory figure is \emph{unified} memory, the
maximum orderable configuration shared with the CPU, which is the
GPU-addressable pool but not dedicated VRAM.

\textbf{Interpretation.} These are \emph{reported nameplate peak} quantities,
not utilized FLOPs: they measure the capacity a paper had access to and chose to
report, an upper bound on work performed. Every capacity figure carries this
caveat.

\textbf{Device counts.} \texttt{device\_count} is stated for $42.9\%$ (S2ORC)
and $46.0\%$ (arXiv) of confirmed mentions, of which $358{,}620$ and $387{,}189$
respectively are explicit rather than inferred. Extraction junk reaches
implausible magnitudes (the raw maxima are $1.5\times 10^8$ and $3\times 10^8$),
so counts are winsorized at the within-year $99.9$th percentile and a hard
ceiling of $65{,}536$ devices. Under the primary \texttt{floor1} policy an
unstated count is imputed as one device, which makes $F_i$ and $M_i$ floors;
an \texttt{explicit}-only variant that drops unstated counts is reported as
robustness. After winsorization the per-paper maximum device count has median 1,
90th percentile 8, and 99th percentile 64 (S2ORC) / 128 (arXiv).

\textbf{Specification coverage.} A paper receives a null on an axis, plus a
\texttt{spec\_missing} flag, when no used model has a specification on that axis;
values are never zero-filled and such papers are excluded from that axis only.
Capacity is computed for $218{,}915$ S2ORC and $201{,}674$ arXiv papers, with
axis coverage FP32 $81.4\%$ / $82.3\%$, FP64 $81.0\%$ / $82.2\%$, FP16 tensor
$57.2\%$ / $64.0\%$, and memory $86.4\%$ / $87.5\%$. Coverage is published by
year alongside every capacity figure, because a coverage trend would otherwise
be indistinguishable from a capacity trend. Tensor capability is treated as a
registry-derived capability flag (a model has a dense tensor throughput figure
or it does not) rather than an inferred property of the text.

\subsection{Denominators and population}

Two denominators are used, and they answer different questions. \emph{Diffusion}
shares are taken over all full-text papers in a field-year --- the inventory
table, which by construction contains papers with no accelerator mention at all.
\emph{Composition} shares (vendor mix, model mix, vintage) are taken over
accelerator-using papers in a field-year. Since corpus size grows by an order of
magnitude over the study period, no count-based series is interpreted without
the corresponding denominator figure.

Publication year for S2ORC comes from the Semantic Scholar metadata snapshot and
for arXiv from the identifier. Field for S2ORC is the first entry of
\texttt{s2fieldsofstudy} ($89.6\%$ of in-window papers have a field label; the
remainder are plotted as a separate ``no field label'' stratum rather than
dropped). Field for arXiv is the archive-level group of the paper's primary
category, from a scraped, versioned copy of arXiv's own taxonomy (8 groups, 155
categories, 18 legacy archive renames), giving $100\%$ coverage; the
fine-grained \texttt{primary\_category} is carried alongside for breakdowns that
have no S2ORC analogue. A test fails if any archive, current or legacy, is
unmapped, since an unmapped archive would silently drop papers out of every
field figure.

The analysis window is $2005$--$2025$ for S2ORC and $2005$--$2026$ for arXiv.
The asymmetry is substantive, not cosmetic: arXiv has no ingestion lag --- a
paper is in the source tars the month it is submitted --- whereas S2ORC's 2026
holds under a tenth of a normal year's papers and would be a pure artefact.
arXiv 2026 is nonetheless a partial calendar year, and is marked as such on
count-based figures; shares are unaffected while numerator and denominator come
from the same partial year.

Citation outcomes are attached only to S2ORC, from a precomputed
outcomes table keyed on Semantic Scholar identifiers, supplying cumulative
citations $k$ years post-publication ($i_k$) and the CD disruption index
($cd_k$). Papers published after the last year with a complete window are
flagged rather than dropped. The arXiv pipeline deliberately maps no arXiv
identifier onto Semantic Scholar, so the arXiv corpus carries no citation
analyses at all.

\subsection{Validity checks and reproducibility}

\textbf{Shared-core regression.} The property that makes the two corpora
comparable at all --- that they are measured by the same code --- is enforced by
tests, not by convention: the same prose fed through an S2ORC record and through
the \LaTeX{} reader must yield identical matched models, candidate counts, and
match offsets that index the stored passage text.

\textbf{Pre-release mentions.} A paper cannot legitimately use a chip released
after it was published, so the pre-release rate is a direct measurement-validity
probe. It is $0.056\%$ of $271{,}380$ S2ORC used-model observations and
$0.027\%$ of $305{,}646$ arXiv observations. Identified causes are metadata
year errors (S2ORC), preprint/version date mismatches, and a small number of
registry release dates off by a year. These observations are flagged, not
silently dropped, and are excluded when reading adoption lags.

\textbf{Converter drift.} Conversion success, encoding, and bibliography
truncation rates are recorded per paper and plotted by year, because a text
pipeline that degrades over eras would manufacture exactly the upward trend this
paper measures.

\textbf{Bibliography leakage.} Candidate passages matching the
reference-shape filter are required to be a negligible share of passages in both
corpora.

\textbf{Provenance.} Every output row carries the registry version, prompt
version, model name and code commit that produced it, and every table is
namespaced by that triple, so no analysis can mix extraction variants.
Recomputation is staged: registry specification changes re-run capacity and
downstream tables only; registry alias changes additionally re-run
canonicalization; only a change to the matcher or the prompt requires re-running
the scan or the LLM.

\textbf{Known limitations.} Reporting is voluntary and unstandardized, so all
estimates are lower bounds on actual accelerator use; capacity is nameplate peak
rather than realized utilization; arXiv any-mention prevalence is biased upward
by raw-source matching; consumer-hardware papers are systematically less likely
to report device counts than cluster papers; and the S2ORC corpus inherits the
field-labelling and text-extraction errors of its own pipeline. We release the
registry, prompts, and mention-level derived records keyed by paper identifier;
full text is not redistributed for either corpus.
